\chapter{Operator Kompak dan Teorema Spektral}
\label{ch:b3:spectral}

Pendiagonalan adalah permata mahkota aljabar linear berdimensi
berhingga: sebab matriks setangkup mempunyai basis ortonormal berisi
vektor eigen. Dalam dimensi tak berhingga hal ini gagal bagi \omterm{def:b3:spectral:selfadjoint}{operator
swa-adjoin} terbatas pada umumnya --- sebab perkalian dengan $x$ pada
$L^2(\intcc01)$ sama sekali \emph{tak} bernilai eigen
(\cref{exo:b3:spectral:6}) --- tetapi ia bertahan, dalam bentuk
yang \omterm{prop:b3:galois:perfect}{sempurna}, bagi operator yang \emph{hampir berdimensi
berhingga}: yakni yang \omterm{def:b3:spectral:compact}{kompak}. Teorema spektral bagi
\omterm{def:b3:spectral:selfadjoint}{operator swa-adjoin} \omterm{def:b3:spectral:compact}{kompak} merupakan teorema yang paling banyak dipakai
dalam analisis fungsional terapan: ia mendiagonalkan persamaan
integral, menggerakkan \omterm{thm:b3:spectral:fredholm}{alternatif Fredholm}, dan (pada soal akhir
pekannya) menyelesaikan dawai yang bergetar, sambil menghasilkan basis sinus
analisis Fourier dari teori operator belaka --- dengan $\zeta(2) =
\frac{\pi^2}6$ milik Euler yang jatuh dari sebuah rumus trace sebagai
hadiah perpisahan. Sepanjang bab ini, $H$ adalah \omterm{def:b3:hilbert:inner}{ruang Hilbert} atas $\C$ (atau
$\R$; pernyataannya menyesuaikan), dan operatornya terbatas.

\section{Operator kompak}

\begin{definition}\label{def:b3:spectral:compact}
Sebuah $T \in \mathcal L(E, F)$ (dengan $E, F$ Banach) disebut
\emph{kompak}\index{operator kompak} jika peta $T(B)$ atas
bola satuannya bersifat kompak relatif di $F$ --- setara dengan itu,
setiap barisan terbatas $(x_n)$ mempunyai subbarisan yang
$(Tx_{n_k})$-nya konvergen. Operator rank berhingga bersifat kompak
(sebab himpunan terbatas dalam dimensi berhingga); sedangkan identitas ruang
berdimensi tak berhingga tak pernah demikian (menurut teorema Riesz, Tahun ke-2).
\end{definition}

\begin{proposition}\label{prop:b3:spectral:ideal}
\omterm{def:b3:spectral:compact}{Operator kompaknya} $\mathcal K(E, F)$ membentuk subruang tertutup
$\mathcal L(E, F)$, sekaligus \omterm{def:b3:rings:ideal}{ideal} dua sisi: yakni $S$ \omterm{def:b3:spectral:compact}{kompak}
$\Rightarrow$ $AS$ dan $SB$ \omterm{def:b3:spectral:compact}{kompak} untuk $A, B$ yang terbatas.
Lebih lanjut, pada \omterm{def:b3:hilbert:inner}{ruang Hilbert}, setiap \omterm{def:b3:spectral:compact}{operator kompak} merupakan
limit \omterm{def:b3:banach:operator}{norma operator} rank berhingga.
\end{proposition}

\begin{proof}
Subruangnya: jelas dari pencirian barisannya. \omterm{def:b3:rings:ideal}{Idealnya}:
sebab pemetaan terbatas mengirim barisan konvergen ke barisan konvergen dan
yang terbatas ke yang terbatas. Ketertutupannya: misalkan $T_n \to T$ dengan $T_n$
\omterm{def:b3:spectral:compact}{kompak}, dan $(x_k)$ terbatas oleh $1$; maka pemetikan diagonal
membuat $(T_nx_{k_j})_j$ konvergen untuk setiap $n$; lalu
$(Tx_{k_j})$ bersifat Cauchy, sebab
\[
\norm{Tx_{k_j} - Tx_{k_l}} \leq 2\vertiii{T - T_n} +
\norm{T_nx_{k_j} - T_nx_{k_l}} ,
\]
dengan memilih $n$ lebih dahulu lalu indeksnya. Untuk penghampiran pada \omterm{def:b3:hilbert:inner}{ruang
Hilbertnya}: misalkan $T$ \omterm{def:b3:spectral:compact}{kompak} dan $K = \overline{T(B)}$ \omterm{def:b3:spectral:compact}{kompak};
lalu diberikan $\varepsilon$, selimuti $K$ oleh berhingga banyak bola
$B(y_i, \varepsilon)$ dan misalkan $P$ proyeksi ortogonal
ke $V = \operatorname{Vect}(y_1, \dots, y_m)$ (yang tertutup sebab
berdimensi berhingga). Maka $PT$ mempunyai rank berhingga, dan untuk
$\norm x \leq 1$: dengan memilih $y_i$ yang $\norm{Tx - y_i} <
\varepsilon$,
\[
\norm{Tx - PTx} \leq \norm{Tx - y_i} + \norm{P(y_i - Tx)}
\leq 2\varepsilon
\]
(sebab $y_i = Py_i$ dan $\vertiii P \leq 1$): sehingga $\vertiii{T - PT} \leq
2\varepsilon$.
\end{proof}

\begin{example}\label{ex:b3:spectral:examples}
(a) Operator diagonal pada $\ell^2$: $T(x_n) = (\lambda_nx_n)$
bersifat \omterm{def:b3:spectral:compact}{kompak} jika dan hanya jika $\lambda_n \to 0$
(\cref{exo:b3:spectral:2}).
(b) Operator kernel pada $\mathcal C(\intcc01)$: \omterm{def:b3:spectral:compact}{kompak} menurut
Ascoli (\cref{exo:b3:complete:7}).
(c) \emph{Operator Hilbert--Schmidt}\index{operator
Hilbert--Schmidt}: untuk $k \in L^2(\intcc01^2)$,
\[
(T_kf)(x) = \int_0^1k(x, y)\,f(y)\,\dd y
\]
mendefinisikan \omterm{def:b3:spectral:compact}{operator kompak} pada $L^2(\intcc01)$ dengan
$\vertiii{T_k} \leq \norm k_{L^2}$
(\cref{exo:b3:spectral:4}: sebab memancung uraian basis
$k$ menampilkan $T_k$ sebagai limit operator rank berhingga).
\end{example}

\section{Operator swa-adjoin}

\begin{definition}\label{def:b3:spectral:selfadjoint}
Sebuah $T \in \mathcal L(H)$ disebut \emph{swa-adjoin}\index{operator
swa-adjoin} jika $T = T^*$ (\cref{exo:b3:hilbert:8}), yakni
$\langle Tx, y\rangle = \langle x, Ty\rangle$ untuk setiap $x, y$.
Maka $\langle x, Tx\rangle \in \R$ untuk setiap $x$ (sebab ia sama dengan
konjugatnya).
\end{definition}

\begin{proposition}\label{prop:b3:spectral:sanorm}
Untuk $T$ yang \omterm{def:b3:spectral:selfadjoint}{swa-adjoin}:\index{jari-jari numerik}
\[
\vertiii T = \sup_{\norm x \leq 1}\ \abs{\langle x,
Tx\rangle} .
\]
Nilai eigen $T$ bersifat real, dan vektor eigen bagi nilai eigen
yang berbeda saling ortogonal.
\end{proposition}

\begin{proof}
Misalkan $M$ supremumnya; maka $M \leq \vertiii T$ menurut
Cauchy--Schwarz. Sebaliknya, identitas bertipe polarisasi
\[
\langle x{+}y, T(x{+}y)\rangle - \langle x{-}y,
T(x{-}y)\rangle = 4\operatorname{Re}\langle y, Tx\rangle
\]
(uraikanlah; sebab suku silangnya $\langle y, Tx\rangle + \langle x,
Ty\rangle = 2\operatorname{Re}\langle y, Tx\rangle$ menurut
\omterm{def:b3:spectral:selfadjoint}{keswa-adjoinannya}) memberikan, bersama hukum jajaran genjangnya,
\[
4\operatorname{Re}\langle y, Tx\rangle \leq
M\bigl(\norm{x{+}y}^2 + \norm{x{-}y}^2\bigr) = 2M\bigl(\norm
x^2 + \norm y^2\bigr).
\]
Lalu untuk $\norm x = 1$ dengan $Tx \neq 0$, ambillah $y = Tx/\norm{Tx}$:
sehingga $4\norm{Tx} \leq 4M$. Jadi $\vertiii T \leq M$. Untuk nilai eigennya:
$Tx = \lambda x$ dengan $x \ne 0$ memberikan $\lambda\norm x^2 = \langle
x, Tx\rangle \in \R$. Untuk keortogonalannya: $\lambda\langle x,
y\rangle = \langle Tx, y\rangle = \langle x, Ty\rangle =
\mu\langle x, y\rangle$ dengan $\lambda \neq \mu$ yang real.
\end{proof}

\section{Teorema spektralnya}

\begin{lemma}[Keberadaan sebuah nilai eigen ekstrem]
\label{lem:b3:spectral:existence}
Misalkan $T \neq 0$ \omterm{def:b3:spectral:compact}{kompak} dan \omterm{def:b3:spectral:selfadjoint}{swa-adjoin}. Maka $\vertiii T$
atau $-\vertiii T$ merupakan nilai eigen $T$.
\end{lemma}

\begin{proof}
Menurut \cref{prop:b3:spectral:sanorm}, pilihlah vektor satuan $x_n$
dengan $\langle x_n, Tx_n\rangle \to \mu$, dengan $\abs\mu =
\vertiii T > 0$ (lalu beralihlah ke subbarisan untuk menetapkan tandanya). Maka
\[
\norm{Tx_n - \mu x_n}^2
= \norm{Tx_n}^2 - 2\mu\langle x_n, Tx_n\rangle + \mu^2
\leq \vertiii T^2 - 2\mu\langle x_n, Tx_n\rangle + \mu^2
\longrightarrow 2\mu^2 - 2\mu\cdot\mu = 0 .
\]
Lewat kekompakannya, sebuah subbarisan $Tx_{n_k} \to y$; lalu $\mu
x_{n_k} = Tx_{n_k} - (Tx_{n_k} - \mu x_{n_k}) \to y$, sehingga
$x_{n_k} \to x = y/\mu$, yakni sebuah vektor satuan, dan \omterm{def:b3:topology:continuity}{kekontinuannya} memberikan
$Tx = \mu x$.
\end{proof}

\begin{theorem}[Teorema spektral bagi operator swa-adjoin
kompak]\label{thm:b3:spectral:spectral}
Misalkan $T$ \omterm{def:b3:spectral:selfadjoint}{operator swa-adjoin} \omterm{def:b3:spectral:compact}{kompak} pada \omterm{def:b3:hilbert:inner}{ruang Hilbert}
$H$.\index{teorema spektral (swa-adjoin kompak)}
\begin{enumerate}
  \item $H$ memiliki sistem ortonormal $(e_n)_{n \in N}$
        ($N$ berhingga atau terbilang) berisi vektor eigen $T$, dengan
        nilai eigen $(\lambda_n)$ yang real dan tak nol, sedemikian sehingga
        \[
        Tx = \sum_{n\in N}\lambda_n\,\langle e_n, x\rangle\,
        e_n \qquad (x \in H),
        \]
        dan $H = \ker T \,\oplus^\perp\,
        \overline{\operatorname{Vect}}(e_n : n \in N)$.
  \item Jika $N$ tak berhingga, maka $\lambda_n \to 0$; lalu untuk setiap
        $\delta > 0$ hanya berhingga banyak $n$ yang mempunyai
        $\abs{\lambda_n} \geq \delta$, dan setiap ruang eigen
        $\ker(T - \lambda)$ dengan $\lambda \neq 0$ berdimensi
        berhingga.
  \item Melengkapkan $(e_n)$ dengan sebuah basis ortonormal $\ker T$
        melahirkan, bila $H$ terpisahkan, sebuah basis ortonormal
        $H$ yang terbuat dari vektor eigen: sehingga $T$ terdiagonalkan.
\end{enumerate}
\end{theorem}

\begin{proof}
(2) lebih dahulu. Seandainya tak berhingga banyak vektor eigen ortonormal $x_k$
mempunyai $\abs{\lambda_{(k)}} \geq \delta$: maka $\norm{Tx_k - Tx_l}^2 =
\lambda_{(k)}^2 + \lambda_{(l)}^2 \geq 2\delta^2$
(menurut keortogonalannya dan Pythagoras): sehingga tak ada subbarisan
$(Tx_k)$ yang konvergen, yang bertentangan dengan kekompakan $T$ pada
$(x_k)$ yang terbatas. Ini membatasi, untuk setiap $\delta$, kelipatan
total nilai eigen di luar
$\intoo{-\delta}\delta$ oleh sebuah bilangan berhingga; sehingga keterbilangannya dan $\lambda_n \to 0$
menyusul.

(1) Misalkan $H_0$ rentang tertutup atas \emph{semua} vektor eigen
bernilai eigen tak nol, yang ditata (menurut (2) dan
Gram--Schmidt di dalam setiap ruang eigen berdimensi berhingga, dengan
keortogonalan antarruang eigennya dari
\cref{prop:b3:spectral:sanorm}) menjadi sistem ortonormal
$(e_n)$ bernilai eigen $\lambda_n \neq 0$. Lalu $T$ memetakan $H_0$
ke $H_0$, dan juga $H_0^\perp$ ke $H_0^\perp$: sebab untuk $y
\perp H_0$ dan $e$ sebuah vektor eigen, $\langle e, Ty\rangle =
\langle Te, y\rangle = \lambda\langle e, y\rangle = 0$. Sedangkan
pembatasan $T' = T\restriction_{H_0^\perp}$ bersifat \omterm{def:b3:spectral:selfadjoint}{swa-adjoin} \omterm{def:b3:spectral:compact}{kompak}
pada \omterm{def:b3:hilbert:inner}{ruang Hilbert} $H_0^\perp$; dan jika $T' \neq
0$, maka \cref{lem:b3:spectral:existence} menghasilkan sebuah vektor eigen
$T$ bernilai eigen tak nol di dalam $H_0^\perp$ ---
yang mustahil, sebab vektor semacam itu tinggal di $H_0$. Jadi $T' = 0$:
sehingga $H_0^\perp \subseteq \ker T$. Sebaliknya $\ker T \perp$ setiap
$e_n$ (sebab $\langle e_n, z\rangle = \frac1{\lambda_n}\langle
Te_n, z\rangle = \frac1{\lambda_n}\langle e_n, Tz\rangle =
0$): jadi $\ker T \subseteq H_0^\perp$, sehingga $\ker T =
H_0^\perp$ beserta penguraian ortogonalnya. Untuk uraiannya:
bagi $x = z + \sum_nc_ne_n$ (dengan $z \in \ker T$ dan $c_n = \langle
e_n, x\rangle$; \cref{thm:b3:hilbert:parseval}(1) pada $H_0$),
\omterm{def:b3:topology:continuity}{kekontinuan} $T$ memberikan $Tx = \sum_nc_n\lambda_ne_n$.

(3) Ruang $\ker T$, sebagai subruang tertutup ruang yang terpisahkan, bersifat
terpisahkan: sehingga ia mempunyai basis ortonormal
(\cref{prop:b3:hilbert:gramschmidt}); dan gabungannya merupakan
basis ortonormal $H$ menurut penguraian pada (1).
\end{proof}

\begin{theorem}[Alternatif Fredholm]
\label{thm:b3:spectral:fredholm}
Misalkan $T$ \omterm{def:b3:spectral:selfadjoint}{swa-adjoin} \omterm{def:b3:spectral:compact}{kompak} dan $\lambda \in \R\setminus
\{0\}$.\index{alternatif Fredholm}
\begin{enumerate}
  \item Jika $\lambda$ bukan nilai eigen, maka $T - \lambda I$
        bersifat bijektif dengan balikan terbatas: sehingga untuk setiap $f$,
        persamaan $Tx - \lambda x = f$ mempunyai tepat satu
        solusi, yang bergantung \omterm{def:b3:topology:continuity}{kontinu} pada $f$.
  \item Jika $\lambda$ merupakan nilai eigen, maka $Tx - \lambda x = f$
        \omterm{def:b3:groups:derived}{terselesaikan} jika dan hanya jika $f \perp \ker(T - \lambda I)$, dan
        solusinya tunggal sampai kernel (berdimensi berhingga)
        itu.
\end{enumerate}
\end{theorem}

\begin{proof}
Uraikan $x = z + \sum c_ne_n$ dan $f = w + \sum d_ne_n$
sepanjang \cref{thm:b3:spectral:spectral} (dengan $z, w \in \ker T$). Maka
persamaannya berbunyi
\[
-\lambda z = w, \qquad (\lambda_n - \lambda)\,c_n = d_n\
(n \in N).
\]
(1) Bila $\lambda \notin \{\lambda_n\}\cup\{0\}$: menurut (2)
teorema spektralnya, $\inf_n\abs{\lambda_n - \lambda} = \delta >
0$ (sebab nilai eigennya menumpuk hanya di $0 \neq \lambda$). Lalu selesaikan:
$z = -w/\lambda$ dan $c_n = d_n/(\lambda_n - \lambda)$, dengan
$\sum\abs{c_n}^2 \leq \delta^{-2}\sum\abs{d_n}^2$: yakni satu-satunya
solusi dengan $\norm x \leq C\norm f$.
(2) Bila $\lambda = \lambda_{n}$ untuk $n$ pada sebuah himpunan berhingga $F$:
\omterm{def:b3:groups:derived}{keterselesaian} $(\lambda_n - \lambda)c_n = d_n$ untuk $n \in F$
menuntut $d_n = 0$, yakni $f \perp e_n$ (bila $n \in F$), yakni $f
\perp \ker(T - \lambda I)$; sedangkan $c_n$ dengan $n \in F$ lalu
bebas.
\end{proof}

\begin{example}\label{ex:b3:spectral:minxy}
Pada $L^2(\intcc01)$, misalkan $Tf(x) = \int_0^1\min(x, y)f(y)\dd y$:
yakni \omterm{ex:b3:spectral:examples}{operator Hilbert--Schmidt} berkernel setangkup real:
jadi \omterm{def:b3:spectral:compact}{kompak} dan \omterm{def:b3:spectral:selfadjoint}{swa-adjoin}. Untuk menyelesaikan $Tf = \lambda f$:
hubungan $\bigl(Tf\bigr)(x) = \int_0^xyf(y)\dd y +
x\int_x^1f(y)\dd y$ menunjukkan bahwa $u = Tf$ memenuhi $u'' = -f$ (lewat dua
penurunan, yang sah untuk $f$ yang \omterm{def:b3:topology:continuity}{kontinu}, dan $Tf$ bersifat
\omterm{def:b3:topology:continuity}{kontinu} untuk $f \in L^2$: lewat kekonvergenan terdominasi), dengan $u(0)
= 0$ dan $u'(1) = 0$. Jadi fungsi eigennya menyelesaikan $\lambda u'' =
-u$, $u(0) = 0$, $u'(1) = 0$:
\[
u_n(x) = \sin\Bigl(\bigl(n + \tfrac12\bigr)\pi x\Bigr),
\qquad
\lambda_n = \frac{1}{\bigl(n + \frac12\bigr)^2\pi^2}
\quad (n \geq 0),
\]
dan teorema spektralnya menyatakan --- tanpa teori Fourier apa pun ---
bahwa sinus ini membentuk basis ortonormal $L^2(\intcc01)$
setelah dinormalkan (sebab kernel $T$ adalah $0$: karena $Tf = 0$
memaksa, lewat kedua penurunannya, $f = 0$ hampir di mana-mana). Sedangkan
soal akhir pekannya menjalankan lingkaran gagasan yang sama untuk
dawai yang bergetar lalu memetik $\zeta(2)$ dari tracenya.
\end{example}

\begin{method}\label{met:b3:spectral:method}
Diberikan sebuah persamaan integral atau diferensial: (1) tuangkanlah kembali sebagai
$(I - \lambda K)u = f$ atau $Ku = \lambda u$ dengan $K$ sebuah
operator integral; (2) periksalah $K$ \omterm{def:b3:spectral:compact}{kompak} (lewat kernel Hilbert--Schmidt
atau Ascoli) dan, bila mungkin, \omterm{def:b3:spectral:selfadjoint}{swa-adjoin} (lewat kernel real
yang setangkup); (3) diagonalkanlah dengan teorema spektralnya atau
panggillah \omterm{thm:b3:spectral:fredholm}{alternatif Fredholm} untuk \omterm{def:b3:groups:derived}{keterselesaiannya}; (4) lalu bacalah
keberadaan, ketunggalan, kestabilan, dan rumus deret bagi
solusinya pada basis eigennya. Operator diferensial bersifat tak
terbatas, tetapi \emph{balikannya} (yakni operator Green) bersifat
\omterm{def:b3:spectral:compact}{kompak}: jadi balikkanlah selalu lebih dahulu.
\end{method}

\section{Latihan}

\begin{exercise}[$\star$]\label{exo:b3:spectral:1}
(a) Tunjukkan bahwa operator terbatas yang jangkauannya berdimensi berhingga
bersifat \omterm{def:b3:spectral:compact}{kompak}.
(b) Tunjukkan bahwa identitas sebuah ruang bernorma bersifat \omterm{def:b3:spectral:compact}{kompak} jika dan hanya jika
dimensinya berhingga (menurut Riesz, Tahun ke-2). Lalu turunkan bahwa \omterm{def:b3:spectral:compact}{operator
kompak} pada ruang berdimensi tak berhingga tak pernah terbalikkan
dengan balikan yang terbatas.
\end{exercise}

\begin{exercise}[$\star$]\label{exo:b3:spectral:2}
Misalkan $T(x_1, x_2, \dots) = (\lambda_1x_1, \lambda_2x_2, \dots)$
pada $\ell^2$, dengan $(\lambda_n)$ yang terbatas.
(a) Tunjukkan $\vertiii T = \sup\abs{\lambda_n}$.
(b) Tunjukkan bahwa $T$ \omterm{def:b3:spectral:compact}{kompak} jika dan hanya jika $\lambda_n \to 0$.
\emph{(Untuk $\Leftarrow$, pancunglah; sedangkan untuk $\Rightarrow$, ujilah pada
$(e_n)$.)}
(c) Kapankah $T$ bersifat \omterm{def:b3:spectral:selfadjoint}{swa-adjoin}? Lalu periksalah teorema spektralnya lewat
pengamatan pada kasus itu.
\end{exercise}

\begin{exercise}[$\star\star$]\label{exo:b3:spectral:3}
Berikan rincian sifat \omterm{def:b3:rings:ideal}{idealnya}
(\cref{prop:b3:spectral:ideal}): bahwa jika $S$ \omterm{def:b3:spectral:compact}{kompak} dan $A, B$
terbatas, maka $ASB$ \omterm{def:b3:spectral:compact}{kompak}. Lalu turunkan bahwa jika $ST = TS = I$
untuk suatu $S$ yang terbatas, dan $\dim H = \infty$, maka $T$ tak
\omterm{def:b3:spectral:compact}{kompak} --- lalu damaikanlah dengan \cref{exo:b3:spectral:1}(b).
\end{exercise}

\begin{exercise}[$\star\star$]\label{exo:b3:spectral:4}
(Hilbert--Schmidt) Misalkan $k \in L^2(\intcc01^2)$ dan $(e_n)$ sebuah
\omterm{def:b3:hilbert:onb}{basis Hilbert} $L^2(\intcc01)$.
(a) Tunjukkan bahwa $\vertiii{T_k} \leq \norm k_{L^2}$
\emph{(lewat Cauchy--Schwarz pada variabel-$y$, lalu Tonelli)}.
(b) Uraikan $k(x,y) = \sum_{m,n}c_{mn}e_m(x)\overline{e_n(y)}$
di $L^2$ pada bujur sangkarnya (benarkanlah bahwa hasil kalinya membentuk
\omterm{def:b3:hilbert:onb}{basis Hilbert} di sana), lalu tunjukkan bahwa memancung jumlahnya memberikan
operator rank berhingga yang konvergen ke $T_k$ dalam \omterm{def:b3:banach:operator}{norma operator}:
sehingga $T_k$ \omterm{def:b3:spectral:compact}{kompak}.
\end{exercise}

\begin{exercise}[$\star\star$]\label{exo:b3:spectral:5}
Misalkan $T$ \omterm{def:b3:spectral:selfadjoint}{swa-adjoin} dengan $\langle x, Tx\rangle \geq 0$ untuk
setiap $x$ (yakni operator \emph{positif}).
(a) Tunjukkan bahwa nilai eigennya $\geq 0$ dan bahwa $\vertiii T =
\sup_{\norm x\leq1}\langle x, Tx\rangle$.
(b) Buktikan ketaksamaan Cauchy--Schwarz yang diperumum
$\abs{\langle x, Ty\rangle}^2 \leq \langle x, Tx\rangle\langle
y, Ty\rangle$.
\end{exercise}

\begin{exercise}[$\star\star$]\label{exo:b3:spectral:6}
Pada $L^2(\intcc01)$, misalkan $(Mf)(x) = x\,f(x)$.
(a) Tunjukkan bahwa $M$ terbatas, \omterm{def:b3:spectral:selfadjoint}{swa-adjoin}, dengan $\vertiii M =
1$, tetapi sama sekali \emph{tak} bernilai eigen.
(b) Tunjukkan bahwa $M$ tak \omterm{def:b3:spectral:compact}{kompak} \emph{(tampilkan sebuah barisan terbatas
yang petanya tak bersubbarisan konvergen, misalnya indikator ternormalkan
selang yang menyusut di dekat $1$ --- atau panggillah
teorema spektralnya)}.
(c) Di manakah bukti \cref{lem:b3:spectral:existence} patah
bagi $M$?
\end{exercise}

\begin{exercise}[$\star\star$]\label{exo:b3:spectral:7}
(Volterra) Pada $L^2(\intcc01)$, misalkan $Vf(x) = \int_0^xf(y)\dd
y$.
(a) Tunjukkan bahwa $V$ \omterm{def:b3:spectral:compact}{kompak} (yakni Hilbert--Schmidt berkernel $\mathbf
1_{y < x}$) tetapi tak \omterm{def:b3:spectral:selfadjoint}{swa-adjoin}; lalu hitunglah $V^*$.
(b) Tunjukkan bahwa $V$ tak bernilai eigen tak nol. \emph{(Dari $Vf =
\lambda f$: $f$ mempunyai wakil yang \omterm{def:b3:topology:continuity}{kontinu}, lalu bersifat
$\mathcal C^1$, dan menyelesaikan $\lambda f' = f$ dengan $f(0) = 0$.)}
(c) Simpulkan bahwa kekompakan saja sama sekali tak melahirkan vektor eigen:
sehingga \omterm{def:b3:spectral:selfadjoint}{keswa-adjoinan} pada \cref{thm:b3:spectral:spectral} bersifat
mendasar.
\end{exercise}

\begin{exercise}[$\star\star\star$]\label{exo:b3:spectral:8}
(Courant--Fischer) Misalkan $T$ \omterm{def:b3:spectral:compact}{kompak}, \omterm{def:b3:spectral:selfadjoint}{swa-adjoin}, dan
\emph{positif}, bernilai eigen $\mu_1 \geq \mu_2 \geq \dots
> 0$ (yang berulang menurut kelipatannya, dengan vektor eigen $e_1, e_2,
\dots$). Tunjukkan bahwa:
\[
\mu_{k} = \max_{\substack{V \subseteq H \\ \dim V = k}}\
\min_{\substack{x \in V\\ \norm x = 1}}\ \langle x, Tx\rangle
= \min_{\substack{W \subseteq H\\ \operatorname{codim}W =
k-1}}\ \max_{\substack{x\in W\\ \norm x = 1}}\ \langle x,
Tx\rangle .
\]
\emph{(Ujilah $V = \operatorname{Vect}(e_1,\dots,e_k)$; sedangkan untuk
batas atasnya iriskan sembarang $V$ dengan
ruang bertipe $\operatorname{Vect}(e_k, e_{k+1}, \dots)$:
sebab pencacahan dimensinya memaksa irisan yang tak nol.)} Lalu turunkan
bahwa nilai eigennya bergantung monoton pada $T$ (yakni $T \leq S
\Rightarrow \mu_k(T) \leq \mu_k(S)$).
\end{exercise}

\begin{exercise}[$\star\star$]\label{exo:b3:spectral:9}
Dengan memakai \cref{thm:b3:spectral:fredholm} untuk $Tf(x) =
\int_0^1\min(x,y)f(y)\dd y$ (\cref{ex:b3:spectral:minxy}):
untuk $\lambda \in \R$ yang mana persamaan integral
\[
f(x) - \lambda\int_0^1\min(x,y)\,f(y)\,\dd y = g(x)
\]
mempunyai tepat satu solusi $f \in L^2$ untuk setiap $g \in L^2$? Apa
yang terjadi pada nilai kecualiannya?
\end{exercise}

\begin{exercise}[$\star\star\star$]\label{exo:b3:spectral:10}
Misalkan $S$ pergeseran pada $\ell^2$
(\cref{exo:b3:banach:1}).
(a) Tunjukkan bahwa $S$ tak bernilai eigen, sedangkan setiap $\lambda$
dengan $\abs\lambda < 1$ merupakan nilai eigen $S^*$ (carilah
vektor eigennya secara eksplisit, yakni barisan geometri).
(b) Baik $S$ maupun $S^*$ tak \omterm{def:b3:spectral:compact}{kompak}: periksalah lewat
pengujian bergaya \cref{exo:b3:spectral:2} pada $(e_n)$.
(c) Berkomentarlah: bahwa bagi operator yang tak \omterm{def:b3:spectral:selfadjoint}{swa-adjoin} dan tak \omterm{def:b3:spectral:compact}{kompak},
lanskap nilai eigennya dapat berupa apa saja dari kosong sampai sebuah
cakram penuh --- sedangkan gagasan yang bertahan adalah \emph{spektrumnya},
yang dipelajari pada kuliah berikutnya.
\end{exercise}

\begin{exercise}[$\star\star$]\label{exo:b3:spectral:11}
(Akar kuadrat) Misalkan $T$ \omterm{def:b3:spectral:compact}{kompak}, \omterm{def:b3:spectral:selfadjoint}{swa-adjoin}, dan positif
(yakni $\langle x, Tx\rangle \geq 0$) pada \omterm{def:b3:hilbert:inner}{ruang Hilbert} $H$, dengan
penguraian spektral $Tx = \sum_n\mu_n\langle e_n,
x\rangle e_n$ (dengan $\mu_n > 0$).
(a) Definisikan $Sx = \sum_n\sqrt{\mu_n}\,\langle e_n, x\rangle
e_n$; lalu tunjukkan bahwa $S$ \omterm{def:b3:spectral:compact}{kompak}, \omterm{def:b3:spectral:selfadjoint}{swa-adjoin}, positif, dengan $S^2
= T$.
(b) Buktikan ketunggalannya: bahwa sembarang $R$ yang \omterm{def:b3:spectral:compact}{kompak}, positif, dan \omterm{def:b3:spectral:selfadjoint}{swa-adjoin}
dengan $R^2 = T$ mengawetkan ruang eigen $T$
\emph{(sebab $RT = R^3 = TR$: sehingga $R$ komutatif dengan $T$, jadi
$R(\ker(T - \mu)) \subseteq \ker(T - \mu)$)}, dan pada
$\ker(T - \mu)$, $R$ merupakan operator positif yang berkuadrat
$\mu\,\mathrm{id}$ pada ruang berdimensi berhingga:
diagonalkanlah ia di sana lalu simpulkan $R =
\sqrt\mu\,\mathrm{id}$ pada setiap ruang eigennya, sehingga $R = S$.
(c) Hitunglah $\sqrt G$ untuk operator dawai $G$ pada
\cref{pb:b3:spectral:1}: kernel mana yang bernilai eigen
$\frac1{n\pi}$ pada basis sinusnya? (Nyatakanlah $\sqrt G$ sebagai
limit-$L^2$ kernelnya; tanpa perlu bentuk tertutup.)
\end{exercise}

\begin{exercise}[$\star\star\star$]\label{exo:b3:spectral:12}
(Penguraian nilai singular) Misalkan $T \in \mathcal L(H)$
\emph{\omterm{def:b3:spectral:compact}{kompak}}, yang tak harus \omterm{def:b3:spectral:selfadjoint}{swa-adjoin}.
(a) Tunjukkan bahwa $T^*T$ \omterm{def:b3:spectral:compact}{kompak}, \omterm{def:b3:spectral:selfadjoint}{swa-adjoin}, dan positif; lalu misalkan
$(e_n)$ keluarga ortonormal berisi vektor eigen dengan
$T^*Te_n = s_n^2e_n$ dan $s_n > 0$ (yakni \emph{nilai singularnya}),
yang dilengkapkan oleh $\ker(T^*T) = \ker T$ (buktikanlah kesamaan ini).
(b) Tetapkan $f_n = \frac{Te_n}{s_n}$; lalu tunjukkan bahwa $(f_n)$
ortonormal, dan tegakkanlah \emph{penguraian nilai singularnya}:
\[
Tx = \sum_n s_n\,\langle e_n, x\rangle\,f_n
\qquad (x \in H),
\]
dengan kekonvergenan di $H$.
(c) Turunkan: bahwa $\vertiii T = \max_ns_n$; bahwa $T$ merupakan limit \omterm{def:b3:banach:operator}{norma
operator} rank berhingga (yang membuktikan kembali
konvers \cref{prop:b3:spectral:ideal} bagi \omterm{def:b3:hilbert:inner}{ruang
Hilbert}); dan untuk operator Volterra $V$ pada
\cref{exo:b3:spectral:7}, yang tak bernilai eigen, terangkanlah
mengapa penguraian nilai singularnya tetap ada dan apa saja bahannya
(kenalilah $V^*V$ sebagai operator kernel bertipe dawai ---
sedangkan menghitung nilai eigennya secara eksplisit adalah
wilayah \cref{exo:b3:spectral:9}).
\end{exercise}

\section{Soal: dawai yang bergetar dan
\texorpdfstring{$\zeta(2)$}{zeta(2)}}

\begin{problem}[{Soal akhir pekan --- operator Green,
basis sinus, dan sebuah rumus trace}]\label{pb:b3:spectral:1}
Kita menyelesaikan soal nilai eigen dawai yang bergetar dengan
ujung tetap --- yakni $-u'' = \nu u$, $u(0) = u(1) = 0$ --- lewat
teori operator, memperoleh basis ortonormal sinusnya tanpa
perhitungan Fourier apa pun, lalu menilaikan $\zeta(2)$ dengan membandingkan dua
ungkapan bagi \emph{trace} operator Greennya.
Definisikan, pada $L^2(\intcc01)$,
\[
(Gf)(x) = \int_0^1 g(x,y)\,f(y)\,\dd y,
\qquad
g(x, y) = \min(x,y)\,\bigl(1 - \max(x,y)\bigr).
\]

\medskip
\noindent\textbf{Bagian I --- Operator Greennya.}
\begin{enumerate}
  \item Tunjukkan bahwa $g$ \omterm{def:b3:topology:continuity}{kontinu}, setangkup, dengan $0 \leq g
        \leq \frac14$, dan bahwa $G$ bersifat \omterm{def:b3:spectral:compact}{kompak} dan
        \omterm{def:b3:spectral:selfadjoint}{swa-adjoin} (\cref{ex:b3:spectral:examples}(c)).
  \item Untuk $f$ yang \omterm{def:b3:topology:continuity}{kontinu}, tunjukkan bahwa $u = Gf$ bersifat $\mathcal
        C^2$ dengan
        \[
        -u'' = f, \qquad u(0) = u(1) = 0
        \]
        \emph{(tulislah $u(x) = (1-x)\int_0^xyf(y)\dd y +
        x\int_x^1(1-y)f(y)\dd y$ lalu turunkan dua kali)}.
        Sebaliknya, jika $u \in \mathcal C^2$ dengan $u(0) = u(1)
        = 0$, maka $G(-u'') = u$: sehingga $G$ membalikkan operator
        dawainya.
  \item Tunjukkan $\ker G = \{0\}$ \emph{(sebab jika $Gf = 0$ dengan $f \in
        L^2$: ujilah terhadap $\varphi$ yang \omterm{def:b3:topology:continuity}{kontinu}, pindahkan $G$
        lewat kesetangkupan/Fubini ke $\varphi$, lalu pakailah
        lema fundamental \cref{cor:b3:lp:fundlemma} --- atau
        regularkanlah)}, dan bahwa $G$ merupakan operator positif:
        yakni $\langle f, Gf\rangle \geq 0$. \emph{(Untuk $f$ yang
        \omterm{def:b3:topology:continuity}{kontinu}: $\langle f, Gf\rangle = \int_0^1 (u')^2$ dengan $u
        = Gf$, lewat pengintegralan parsial; lalu simpulkan lewat \omterm{ex:b3:lebesgue:gamma}{kepadatannya}.)}
\end{enumerate}

\medskip
\noindent\textbf{Bagian II --- Pendiagonalannya: basis
sinusnya.}
\begin{enumerate}[resume]
  \item Tunjukkan bahwa fungsi eigen $G$ bernilai eigen
        $\lambda \ne 0$ merupakan, sampai skalarnya, solusi
        $-\lambda u'' = u$, $u(0) = u(1) = 0$ \emph{(sebab sebuah
        fungsi eigen mempunyai wakil yang \omterm{def:b3:topology:continuity}{kontinu} --- karena $Gf$
        \omterm{def:b3:topology:continuity}{kontinu} untuk $f \in L^2$, mengapa? --- sehingga ia
        $\mathcal C^2$ lewat pengangkatan pertanyaan 2)}.
  \item Selesaikan soal nilai batasnya: bahwa nilai eigen
        $G$ adalah $\lambda_n = \frac1{n^2\pi^2}$ (dengan $n \geq 1$),
        dengan fungsi eigen ternormalkan $e_n(x) =
        \sqrt2\,\sin(n\pi x)$; lalu periksalah keortonormalannya lewat
        pengintegralan langsung sebagai uji kewarasan.
  \item Simpulkan dari
        \cref{thm:b3:spectral:spectral} dan pertanyaan 3 bahwa
        $\bigl(\sqrt2\sin(n\pi x)\bigr)_{n\geq1}$ merupakan
        \emph{basis} ortonormal $L^2(\intcc01)$ --- tanpa
        Stone--Weierstrass dan tanpa deret Fourier. Lalu uraikan
        $f(x) = x(1-x)$ pada basis ini dan tuliskan \omterm{thm:b3:hilbert:parseval}{Parseval}
        baginya.
\end{enumerate}

\medskip
\noindent\textbf{Bagian III --- Rumus tracenya dan
$\zeta(2)$.}
\begin{enumerate}[resume]
  \item Buktikan kedua identitas
        \[
        \langle e_n, Ge_n\rangle = \lambda_n
        \quad\text{dan}\quad
        \sum_{n\geq1}\lambda_n = \int_0^1 g(x,x)\,\dd x .
        \]
        Untuk yang kedua (yakni \emph{rumus tracenya}): uraikan
        $g(x, \cdot)$, untuk $x$ yang tetap, pada basis $(e_n)$ ---
        lalu tunjukkan bahwa koefisiennya adalah $\lambda_ne_n(x)$, sehingga
        $g(x, \cdot) = \sum_n\lambda_ne_n(x)\,e_n$ di
        $L^2$. Di sini sinusnya eksplisit: periksalah secara langsung
        bahwa $\sum_n\lambda_ne_n(x)e_n(y)$ konvergen
        \emph{secara seragam} pada bujur sangkarnya (bandingkan dengan $\sum
        \frac2{n^2\pi^2}$), sehingga jumlahnya \omterm{def:b3:topology:continuity}{kontinu} dan, karena
        beruraian-$L^2$ sama pada $y$ untuk setiap $x$, sama dengan
        $g(x,y)$ di mana-mana. Lalu tetapkan $y = x$ dan integralkan suku demi
        suku.
  \item Hitunglah $\int_0^1g(x,x)\dd x = \int_0^1x(1-x)\dd x =
        \frac16$, lalu simpulkan
        \[
        \sum_{n\geq1}\frac{1}{n^2\pi^2} = \frac16,
        \qquad\text{yakni}\qquad
        \boxed{\ \zeta(2) = \frac{\pi^2}{6}\ } :
        \]
        yakni jumlah Euler dari sebuah trace operator.
  \item Turunkanlah $\zeta(2)$ dengan cara ketiga: terapkan \omterm{thm:b3:hilbert:parseval}{Parseval} pada
        basis sinusnya untuk fungsi konstan $\mathbf 1$,
        hitunglah $\sum_{n \text{ gasal}}\frac1{n^2}$, lalu
        simpulkan. Lalu bandingkan mekanismenya: dalam arti apakah
        argumen trace pertanyaan 7--8 merupakan ``\omterm{thm:b3:hilbert:parseval}{Parseval}
        yang diterapkan pada seluruh kernelnya sekaligus''?
\end{enumerate}

\medskip
\noindent\textbf{Bagian IV --- Dawainya bergetar.}
\begin{enumerate}[resume]
  \item (Pemisahan variabel, yang disintesiskan) Untuk $f \in
        L^2$, definisikan
        \[
        u(t, x) = \sum_{n\geq1}\;c_n\,
        \cos(n\pi t)\,\sqrt2\sin(n\pi x),
        \qquad c_n = \langle e_n, f\rangle .
        \]
        Tunjukkan bahwa deretnya konvergen di $L^2(\intcc01)$ untuk setiap
        $t$, bahwa $t\mapsto u(t, \cdot)$ \omterm{def:b3:topology:continuity}{kontinu} ke dalam
        $L^2$, dan bahwa untuk $f$ pada rentang berhingga banyak
        $e_n$ ia menyelesaikan persamaan gelombang $\partial_t^2u =
        \partial_x^2u$ dengan $u(0) = f$, $\partial_tu(0) = 0$,
        dan ujung yang tetap. Nilai eigennya $n^2\pi^2$ adalah kuadrat
        frekuensinya: yakni harmonik dawainya --- lalu terangkanlah
        tafsiran musik \cref{thm:b3:spectral:spectral} dalam satu paragraf.
\end{enumerate}

\medskip
\noindent\textbf{Bagian V --- Buah hasil variasinya: metode
pangkat, kestabilan Weyl, dan sebuah batas ketat bagi $\pi$.} Misalkan
$A$ \omterm{def:b3:spectral:selfadjoint}{operator swa-adjoin} \omterm{def:b3:spectral:compact}{kompak} yang \emph{positif} dengan
nilai eigen $\mu_1 \geq \mu_2 \geq \cdots > 0$ dan
vektor eigen ortonormal $(u_n)$; lalu $R(x) = \frac{\langle x,
Ax\rangle}{\norm x^2}$. Rumus min--maksnya adalah
\cref{exo:b3:spectral:8}; dan di sini kita membelanjakannya.
\begin{enumerate}[resume]
  \item (Metode pangkat) Untuk $x \neq 0$ tulislah $m_p =
        \sum_n\mu_n^p\abs{\langle u_n, x\rangle}^2$. Tunjukkan
        $m_pm_{p+2} \geq m_{p+1}^2$ (lewat Cauchy--Schwarz), lalu turunkan
        rantainya
        \[
        R(x) \;\leq\; \frac{\langle Ax, Ax\rangle}{\langle x,
        Ax\rangle} \;\leq\; R(Ax) \;\leq\; \mu_1,
        \]
        lalu buktikan bahwa jika $\langle u_1, x\rangle \neq 0$,
        maka $R(A^kx) \to \mu_1$: sehingga mengiterasikan operatornya pada
        sembarang vektor generik menghitung nilai eigen puncaknya ---
        yakni metode pangkat analisis numerik, yang tersahihkan.
  \item (Kestabilan Weyl) Untuk $A, B$ yang \omterm{def:b3:spectral:selfadjoint}{swa-adjoin} \omterm{def:b3:spectral:compact}{kompak} dan positif,
        turunkan dari \cref{exo:b3:spectral:8} bahwa
        \[
        \abs{\mu_n(A) - \mu_n(B)} \;\leq\; \vertiii{A - B}
        \qquad\text{untuk setiap } n :
        \]
        sehingga seluruh spektrumnya bersifat Lipschitz-$1$ dalam \omterm{def:b3:banach:operator}{norma
        operatornya} --- jadi nilai eigen sistem setangkup yang besar dapat
        dihitung dari hampirannya dengan galat yang terjamin.
  \item Terapkan batas Rayleigh pada $G$ dengan fungsi uji
        $u(x) = x(1-x)$: selesaikanlah $-w'' = u$ dengan $w(0) =
        w(1) = 0$ untuk memperoleh $Gu = w = \frac{x(1-x)(1 + x -
        x^2)}{12}$, lalu hitunglah
        \[
        \norm u_2^2 = \frac1{30}, \qquad
        \langle u, Gu\rangle = \frac1{12}\Bigl(\frac1{30} +
        \frac1{140}\Bigr) = \frac{17}{5040},
        \qquad R(u) = \frac{17}{168},
        \]
        lalu simpulkan batas \emph{ketatnya} $\frac1{\pi^2}
        = \lambda_1 \geq \frac{17}{168}$, yakni $\pi \leq
        \sqrt{168/17} < 3.1437$.
  \item Satu langkah rantai pertanyaan 11, pada fungsi uji
        yang sama: dengan memakai $\int_0^1(x - x^2)^4\dd x =
        \frac1{630}$, hitunglah
        \[
        \norm{Gu}_2^2 = \frac1{144}\Bigl(\frac1{30} +
        \frac2{140} + \frac1{630}\Bigr) = \frac{31}{90720},
        \qquad
        \frac{\langle Gu, Gu\rangle}{\langle u, Gu\rangle} =
        \frac{31}{306},
        \]
        lalu simpulkan $\pi \leq \sqrt{306/31} < 3.1419$: yakni dua
        integral, empat angka yang benar. (Setiap iterasi berikutnya
        kira-kira mengkuadratkan ketelitiannya: sebab jurang vektor eigennya
        $\lambda_1/\lambda_2 = 4$ menggerakkan kekonvergenan
        geometrinya.)
\end{enumerate}

\medskip
\noindent\textbf{Bagian VI --- Jejak $G^2$, dan
$\zeta(4)$.}
\begin{enumerate}[resume]
  \item Tunjukkan bahwa $\sum_n\lambda_n^2 =
        \iint_{\intcc01^2}g(x,y)^2\,\dd x\,\dd y$
        \emph{(uraikan $g$ pada basis hasil kalinya
        $(e_m(x)e_n(y))_{m,n}$ dari $L^2(\intcc01^2)$ --- yakni
        sebuah \omterm{def:b3:hilbert:onb}{basis Hilbert}, bandingkan \cref{exo:b3:spectral:5} --- lalu
        terapkan \omterm{thm:b3:hilbert:parseval}{Parseval} pada bujur sangkarnya; sedangkan pertanyaan 7 mengenali
        koefisiennya)}.
  \item Hitunglah integral rangkapnya:
        \[
        \iint g^2 = 2\int_0^1(1-x)^2\Bigl(\int_0^x
        y^2\,\dd y\Bigr)\dd x
        = \frac23\int_0^1x^3(1-x)^2\,\dd x = \frac1{90} .
        \]
  \item Simpulkan $\zeta(4) = \dfrac{\pi^4}{90}$; lalu terangkan,
        tanpa perhitungan, bagaimana trace pangkat yang lebih tinggi
        $G^k$ menghasilkan $\zeta(2k) \in \pi^{2k}\,\Q$ untuk setiap
        $k \geq 1$, dan mengapa nilai gasalnya $\zeta(3),
        \zeta(5), \dots$ secara struktural berada di luar
        jangkauan mesin ini.
  \item ($\pi$ dari bawah) Dari $\lambda_1^2 \leq
        \sum_n\lambda_n^2 = \frac1{90}$ turunkan $\pi \geq
        90^{1/4} > 3.080$, lalu rakitlah bersama pertanyaan 14
        putusan dua sisinya
        \[
        3.080 \;<\; \pi \;<\; 3.1419,
        \]
        yang diperoleh sepenuhnya dari aritmetika dawai yang
        bergetar. Sisi mana yang konvergen lebih cepat bila orang memakai
        trace yang lebih tinggi $(\operatorname{tr}G^{2k})^{-1/4k}$,
        dan mengapa?
\end{enumerate}

\medskip
\noindent\textbf{Bagian VII --- Penggerakan dan resonansi.} Tetapkan
$\nu \in \R$ lalu tinjaulah dawai yang digerakkan $-u'' - \nu u =
f$, $u(0) = u(1) = 0$, dengan $f \in L^2$ dan $c_n = \langle
e_n, f\rangle$.
\begin{enumerate}[resume]
  \item Andaikan $\nu \notin \{n^2\pi^2 : n \geq 1\}$. Tunjukkan
        bahwa
        \[
        u = \sum_{n\geq1}\frac{c_n}{n^2\pi^2 - \nu}\,e_n
        \]
        konvergen di $L^2$ dan secara seragam pada $\intcc01$
        \emph{(lewat Cauchy--Schwarz antara $(c_n)$ dan ekor
        $\sum n^{-4}$, dengan $\norm{e_n}_\infty =
        \sqrt2$)}, dan bahwa ia memenuhi $u = Gf + \nu Gu$
        --- yakni bentuk berkoordinat eksplisit \omterm{thm:b3:spectral:fredholm}{alternatif
        Fredholm} (\cref{thm:b3:spectral:fredholm}), beserta
        ketunggalannya.
  \item Andaikan $\nu = m^2\pi^2$. Tunjukkan bahwa $u = Gf + \nu
        Gu$ bersolusi jika dan hanya jika $c_m = 0$, dan tunggal sampai
        penambahan kelipatan $e_m$. Bacaan fisikanya: yakni mendorong
        ayunan tepat pada frekuensinya sendiri.
  \item Untuk $\nu < \pi^2$, tunjukkan bahwa operator solusinya
        $R_\nu\colon f \mapsto u$ terbatas pada $L^2$ dengan
        norma $\frac1{\pi^2 - \nu}$, \omterm{def:b3:spectral:compact}{kompak}, \omterm{def:b3:spectral:selfadjoint}{swa-adjoin},
        dan positif: sehingga seluruh analisis spektralnya berulang,
        tergeser sejauh $\nu$.
  \item (Sintesis) Susunlah kamus soal ini:
        nilai eigen $\leftrightarrow$ kuadrat frekuensi
        (harmonik); trace $\leftrightarrow$ $\zeta(2)$;
        norma Hilbert--Schmidt $\leftrightarrow$ $\zeta(4)$;
        \omterm{thm:b3:spectral:fredholm}{alternatif Fredholm} $\leftrightarrow$ resonansi;
        min--maks $\leftrightarrow$ batas variasi (dengan $\pi <
        3.1437$ dari satu polinomial). Yakni satu operator integral,
        lima bab analisis yang diuangkan.
\end{enumerate}

\medskip
\noindent\textbf{Bagian VIII --- Tiga gema terakhir.}
\begin{enumerate}[resume]
  \item ($\zeta(6)$, secara cuma-cuma) Identitas \omterm{thm:b3:hilbert:parseval}{Parseval}
        pertanyaan 6 memberikan $\sum_{n\text{ gasal}}n^{-6} =
        \frac{\pi^6}{960}$. Pisahkan $\zeta(6)$ menurut $n$ yang gasal dan
        genap lalu simpulkan
        \[
        \zeta(6) = \frac{\pi^6}{945},
        \]
        tanpa integral baru: sebab mesin pertanyaan 17
        (yakni trace $G^3$) akan menghasilkan nilai yang sama
        dengan biaya sebuah kernel teriterasi --- jadi \omterm{thm:b3:hilbert:parseval}{Parseval} pada satu
        fungsi yang terpilih baik adalah jalur yang lebih murah di sini.
  \item (Keadaan dasarnya positif) Misalkan $A$ \omterm{def:b3:spectral:selfadjoint}{operator swa-adjoin}
        \omterm{def:b3:spectral:compact}{kompak} positif pada $L^2(\intcc01)$
        yang diberikan oleh kernel setangkup \omterm{def:b3:topology:continuity}{kontinu} $k > 0$ pada
        $\intoo01^2$, bernilai eigen terbesar $\mu_1$. Tunjukkan bahwa:
        (a) sembarang pemberi maksimum hasil bagi Rayleighnya merupakan
        fungsi eigen-$\mu_1$; (b) jika $u$ salah satunya, maka
        $\langle\abs u, A\abs u\rangle \geq \langle u,
        Au\rangle$, dengan ketaksamaan \emph{sejati} jika $u$ bertanda
        dua-duanya pada himpunan berukuran positif --- sehingga $u$
        bertanda tetap hampir di mana-mana, dan $u = \mu_1^{-1}Au$
        tak pernah lenyap pada $\intoo01$; (c) $\mu_1$ merupakan
        nilai eigen yang \emph{sederhana}. Lalu periksalah setiap klaimnya pada $G$:
        sebab $e_1 = \sqrt2\sin(\pi x) > 0$, dan setiap $e_n$ dengan $n
        \geq 2$, karena ortogonal terhadap $e_1$, haruslah berganti tanda
        (dan memang: sebab ia berakar $n - 1$ kali di \omterm{def:b3:topology:interior}{interiornya}).
  \item (Jarak ke spektrumnya, dan harga
        resonansinya) Untuk $\nu \notin \{n^2\pi^2\}$, tunjukkan bahwa
        operator solusi $R_\nu$ pada pertanyaan 19 bersifat
        terbatas, \omterm{def:b3:spectral:selfadjoint}{swa-adjoin}, \omterm{def:b3:spectral:compact}{kompak}, dengan
        \[
        \vertiii{R_\nu} =
        \frac1{\min_{n\geq1}\,\abs{n^2\pi^2 - \nu}}
        = \frac1{\operatorname{dist}\bigl(\nu,
        \{n^2\pi^2\}\bigr)},
        \]
        dengan normanya tercapai pada ragam yang terdekat. Lalu
        ukurlah ayunan pertanyaan 20: sebab menggerakkannya dengan $f = e_1$
        pada $\nu = (1 - \varepsilon)\pi^2$ menghasilkan $u =
        \frac{1}{\varepsilon\pi^2}\,e_1$, yakni penguatan sebesar
        $\frac1\varepsilon$ atas tanggapan statisnya $Ge_1 =
        \frac1{\pi^2}e_1$ --- sehingga pada satu persen di bawah
        nada dasarnya (yakni $\varepsilon = 10^{-2}$), dawainya
        menjawab seratus kali lebih keras.

\end{enumerate}
\end{problem}
